\documentclass[a4paper,11pt]{article}
\usepackage[T1]{fontenc}
\usepackage[utf8]{inputenc}
\usepackage{lmodern}
\usepackage{amsmath,amsthm,amssymb,amsfonts}
\usepackage{hyperref}

\newtheorem{claim}{Claim}
\newtheorem{definition}{Definition}

\title{Graduate Analysis 1: midsem exam}
\author{Aditya Padiyar}
\date{October 2026}

\begin{document}

\maketitle

\begin{definition}
  A sequence of functions $f_n$ converges almost uniformly to $f$ 
  if for every $\epsilon>0$ there is a set $A_\epsilon$ with $\mu(A_\epsilon)<\epsilon$ such that 
  $f_n$ converges uniformly to $f$ on $(A_\epsilon)^c$
\end{definition}

\begin{definition}
  A sequence $f_n \in L^1$ is \textbf{uniformly integrable} if 
  for all $\epsilon>0$ there is $\delta>$ such that $\int_A |f_n| < \epsilon$ for all sets $A$ with $\mu(A)<\delta$
\end{definition}

\begin{claim}\label{uniform_sum}
  The sum of two uniformly integrable sequences is uniformly integrable
\end{claim}

\begin{claim}
  On a finite measure space $(X,\Sigma,\mu)$, convergence almost everywhere is equivalent to convergence almost uniformly
\end{claim}

\begin{proof}
  Consider the following:
  \begin{itemize}
    \item See Egorov's theorem in Axler for a proof of the forward implication
    \item Let $B \in \Sigma$ denote the points where convergence does not occur. Then $B \subseteq A_\epsilon$ which implies that $\mu(B)=0$
  \end{itemize}
\end{proof}

\begin{claim}
  Suppose $f_n^0,f_n^1,f_n^2,f_0,f_1,f_2$ are functions such that 
  \begin{itemize}
    \item $f_n^i,f_0,f_2 \in L^1$ 
    \item $f_n^0 \leq f_n^1 \leq f_n^2$
    \item $f_n^i \to f_i$
    \item $\int f_n^0 \to \int f_0$ and $\int f_n^2 \to \int f_2$
  \end{itemize}
  
  Then we have 
  \begin{itemize}
    \item $f_1 \in L^1$ 
    \item $\int f_n^1 \to \int f_1$
  \end{itemize}
\end{claim}

\begin{proof}
  Consider the following:
  \begin{itemize}
    \item Taking limits we have $f_0 \leq f_1 \leq f_2$. Therefore $|f_1| \leq |f_0| + |f_2|$ which shows that $f_1 \in L^1$ 
    \item Define the following:
          \begin{itemize}
            \item $g_n=f_n^1-f_n^0$
            \item $h_n=f_n^2-f_n^1$
            \item $g=f_1-f_0$
            \item $h=f_2-f_1$
          \end{itemize}
    \item By Fatou's lemma $\int g \leq \text{liminf} \int g_n$ and $\int h \leq \text{liminf} \int h_n$
    \item We have \begin{align*}
                    \text{limsup} \int g_n = \text{limsup} \int (g_n +h_n -h_n) \\
                    = \int f_2 - \int f_0 - \text{liminf} \int h_n \leq \int g 
                  \end{align*}
    \item Hence $\int g_n \to \int g$
    \item We have $\int f_n^1 = \int g_n + \int f_n^0 \to \int g + \int f_0 = \int f_1$
  \end{itemize}
\end{proof}

\begin{claim}
  Suppose $f:\mathbb{R}^2 \to \mathbb{R}$ is measurable and satisfies:
  
  \begin{itemize}
    \item $[x \to f(x_0,x)] \in L^1$ for all $x_0$
    \item $g(x,y)=\partial_x f(x,y)$ exists everywhere
    \item There is $h \in L^1$ on $\mathbb{R}$ such that $|g(x,y)| \leq h(y)$ 
  \end{itemize}
  
  Let $\phi(x)= \int f(x,y)dy$. Then $\phi'(x)=\int g(x,y)dy$
\end{claim}

\begin{proof}
  Consider the following:
  \begin{itemize}
    \item Fix $x \in \mathbb{R}$ and define $q(t,y)=\frac{f(x+t,y)-f(x,y)}{t}$ for $t \neq 0$
    \item By the mean value theorem $q(t,y)=g(\alpha(t,y),y)$
    \item Pick $t_n \neq 0$ such that $t_n \to 0$
    \item We have $\frac{\phi(x+t_n)-\phi(x)}{t_n} = \int q(t_n,y)dy$
    \item We have $q(t_n,y) \to g(x,y)$ for every $y$
    \item By DCT $\phi'(x)= \int g(x,y)dy$
  \end{itemize}
\end{proof}

\begin{claim}
  Suppose $X$ is a finite measure space.
  If $f_n:X \to \mathbb{R}$ are uniformly integrable and converge pointwise to $f:X \to \mathbb{R}$, then $\int |f_n-f| \to 0$
\end{claim}

\begin{proof}
  Consider the following:
  \begin{itemize}
    \item We may assume $f=0$ by 3.28 of Axler and \ref{uniform_sum}
    \item Let $\alpha > 0$ and $\delta_\alpha$ be such that $\mu(A)<\delta_\alpha$ makes $\int_A |f_n|<\frac{\alpha}{2}$
    \item Choose $A_\alpha$ with $\mu(A)<\delta_\alpha$ such that $f_n$ converges uniformly to $0$ on $(A_\alpha)^c$
    \item Then \[\int |f_n| = \int_{A_\alpha} |f_n| + \int_{(A_\alpha)^c} |f_n| \leq \frac{\alpha}{2} + \int_{(A_\alpha)^c} |f_n| \to \frac{\alpha}{2}\]
  \end{itemize}
\end{proof}



\begin{claim}
  Suppose $X$ is a random variable with finite expectation and $F \subseteq \mathbb{R}$ is closed.
  Suppose that whenever $P(E)>0$ we have $\frac{\mathbb{E}[X1_E]}{P(E)} \in F$. Then $P(X \in F)=1$
\end{claim}

\begin{proof}
  Consider the following:
  \begin{itemize}
    \item Suppose $P(X \notin F)>0$
    \item Then there is a compact interval $[a,b] \subseteq F^c$ such that $P(X \in [a,b])>0$
    \item Let $E$ be the event $\{w: X(w) \in [a,b]\}$
    \item We have $\frac{\mathbb{E}[X1_E]}{P(E)} \in [a,b]$ which is a contradiction
  \end{itemize}
\end{proof}

\begin{claim}
  Define the Cantor set as in 2.74 of Axler. Define $f=\sum_{n=1}^\infty n1_{G_n}$. Then we have
  \begin{itemize}
    \item $f$ is measurable
    \item $\int f = 3$
  \end{itemize}
\end{claim}

\begin{proof}
  Define $f_k = \sum_{n=1}^k n1_{G_n}$. Then $f_k$ converges to $f$ pointwise. 
  This makes $f$ measurable. By the monotone convergence theorem, \[\int f = \sum_{n=1}^\infty n\lambda(G_n)\]
  Since $\lambda(G_n)=\frac{1}{2}(\frac{2}{3})^n$, the result follows
\end{proof}

\begin{claim}
  Suppose $E \subseteq [0,1] \times [0,1]$. If $\lambda\{x:\lambda(E_x)\leq\frac{1}{2}\}=1$ then $\lambda\{y:\lambda(E^y)=1\}\leq \frac{1}{2}$
\end{claim}

\end{document}
